Chapter~\ref{chap:distributions} identified three distribution types.
Sequential and near-sequential keys have small local spread, so exact mode
(Section~\ref{sec:adaptive}) applies directly --- FPR~$= 0$ with no hash needed.
When distribution is uniformed and the key span is large --- keys are spread out, truncation
(Section~\ref{sec:truncation}) is effective: large gaps between keys mean few phantom
collisions and low FPR.

Clustered keys (OSM, geospatial IDs) fit neither case alone: dense regions inside each
cluster behave like near-sequential data, while the sparse gaps between clusters have
large inter-key distances where truncation works well.

The hybrid idea is to split along this structure --- isolate each dense region into its
own exact-mode sub-filter (FPR~$= 0$), and route the sparse keys between clusters to
a truncation fallback.

\begin{figure}[h]
    \centering
    \includegraphics[width=0.85\textwidth]{segmentation}
    \caption{Hybrid segmentation: dense clusters go into exact-mode sub-filters
        (FPR~$= 0$); sparse keys between clusters go to a truncation fallback.}
    \label{fig:segmentation}
\end{figure}

The rest of this chapter presents the segmentation algorithm: 1D DBSCAN
(Section~\ref{sec:hybrid-scan}).

\section{1D DBSCAN Segmentation (Scan-ARE)}
\label{sec:hybrid-scan}

\subsection{1D DBSCAN in $O(n)$}

DBSCAN~\citep{ester1996dbscan} is normally $O(n^2)$ due to neighbor search. On sorted
1D data, the $\delta$-neighborhood of any point is an interval, so the
algorithm reduces to two-pointer sweeps:

\begin{enumerate}
  \item \textbf{Forward sweep:} for each $\text{key}[i]$, advance the left pointer while
    $\text{key}[i] - \text{key}[\text{left}] > \delta$. If the window has
    $\geq \texttt{minPts}$ points, $\text{key}[i]$ is a core point.
  \item \textbf{Reverse sweep:} same logic right-to-left (the forward pass only marks
    the rightmost points of each dense window).
  \item \textbf{Core runs:} contiguous core points with gap $\leq \delta$ form
    cluster cores.
  \item \textbf{Merge:} adjacent cores within $\delta$ of each other are joined.
  \item \textbf{Border expansion:} non-core points within $\delta$ of a core point
    (\emph{border points}) join the nearest cluster.
  \item \textbf{Post-filter:} clusters with fewer than \texttt{minClusterSize} keys
    dissolve back to fallback.
\end{enumerate}

\begin{table}[H]
\centering
\begin{tabular}{@{}lrl@{}}
\toprule
Parameter & Value & Role \\
\midrule
\texttt{minPts}         & 10  & DBSCAN core threshold \\
\texttt{minClusterSize} & 256 & Post-filter: clusters below this dissolve \\
\bottomrule
\end{tabular}
\caption{DBSCAN parameters. Two separate parameters --- unlike classic DBSCAN which
  conflates core density and cluster size.}
\label{tab:dbscan-params}
\end{table}

\subsection{Choosing $\delta$ From Problem Parameters}

\begin{equation}
  \delta = \texttt{minPts} \cdot \frac{\mathcal{L}}{\varepsilon}
  \label{eq:delta}
\end{equation}

From the exact mode analysis (Equation~\ref{eq:exact-threshold}): exact mode triggers
when the average gap $< \mathcal{L}/\varepsilon$. Setting $\delta = \texttt{minPts} \cdot \mathcal{L}/\varepsilon$
ties the neighbourhood radius directly to the core threshold: a window must hold
\texttt{minPts} points, and each consecutive pair can be at most $\mathcal{L}/\varepsilon$ apart.

\begin{table}[H]
\centering
\begin{tabular}{@{}lrrl@{}}
\toprule
Distribution & Typical gap & Gap vs $\delta$ & Result \\
\midrule
Cluster (internal gap ${\sim}100$)  & 100            & $\ll \delta$ & Cluster found \\
Uniform 64-bit (avg gap ${\sim}2^{44}$) & ${\sim}10^{13}$ & $\gg \delta$ & All fallback \\
Sequential (gap $= 1000$)          & 1000           & $\ll \delta$ & One big cluster \\
\bottomrule
\end{tabular}
\caption{Sanity check for $\mathcal{L} = 128$, $\varepsilon = 10^{-3}$,
  $\delta = 1.28 \times 10^6$. No data statistics or percentiles needed --- $\delta$ is derived
  purely from $\mathcal{L}$ and $\varepsilon$.}
\label{tab:eps-sanity}
\end{table}

\subsection{Dual Fallback}

Keys not assigned to any cluster go to a fallback filter. The choice between truncation
and SODA is made automatically:

\begin{enumerate}
  \item Let $\text{spread} = \max(\text{fallback keys}) - \min(\text{fallback keys})$.
    Compute phantom size: $\text{phantom\_size} = \text{spread} / 2^K$.
  \item Compute the 5th-percentile gap ($P_5$) of fallback keys via quickselect.
  \item If $P_5 > \text{phantom\_size}$: \textbf{truncation} (safe, saves
    $\log_2(\mathcal{L})$ BPK).
  \item If $P_5 \leq \text{phantom\_size}$: \textbf{adaptive/SODA} (guaranteed FPR).
\end{enumerate}

\subsection{Query}

Clusters are sorted by $[\min, \max]$. For a query $[a, b]$:
\begin{enumerate}
  \item Binary search over clusters intersecting $[a, b]$ --- $O(\log C)$ where $C$ is
    the number of clusters.
  \item Check each intersecting cluster (typically 0--1).
  \item Check the fallback filter.
  \item Return ``empty'' only if all sub-filters agree.
\end{enumerate}

\subsection{FPR Guarantee}

Each sub-filter independently targets FPR~$\leq \varepsilon$. By a union bound, if a
query touches at most $C$ sub-filters, the overall FPR~$\leq C\varepsilon$. In practice
most queries touch only one sub-filter (either one cluster or just the fallback), so
$C = 1$ and FPR~$\approx \varepsilon$.
